\section{Full enclosures and smooth obstacles}
\label{sec:g-enclosures}

A smooth inner boundary need not minimize the area of its enclosures.
The three lemmas in this section replace it by a full minimizing
frontier, retaining every component and every portion of contact with
the original boundary. A strictly negative inner mean curvature then
makes the minimizing frontier detach from that boundary. These are
geometric results for a fixed smooth metric; they require neither
nonnegative scalar curvature nor a mass hypothesis.

We state the results here for later reference.
Their proofs are given together in Section~\ref{sec:g-enclosure-proofs},
immediately before their use in the passage from harmonic-flat to
general asymptotically flat ends. The local obstacle regularity is the
same input that underlies minimizing hulls in Huisken--Ilmanen's
inverse mean curvature flow \cite{HuiskenIlmanen2001}.

\subsection{The full area and its minimizing enclosure}

Let $(X,g)$ be a smooth connected Riemannian manifold of
dimension $n\ge3$, complete with its nonempty compact smooth boundary
$B$ included.  Assume that $X$ has one coordinate end with compact
complement and that, in its coordinates,
\begin{equation}
 g_{ij}-\delta_{ij}=O_2(r^{-q})\qquad\text{for some }q>0.
 \label{eq:g-end}
\end{equation}
Only the eventual strict mean convexity of its coordinate spheres is
needed for the three enclosure lemmas below. In particular, scalar
curvature and ADM mass play no role in those lemmas.

A \emph{full smooth enclosing cut} is the entire intrinsic manifold
boundary $\partial_{\mathrm{man}}D$ of a connected smooth domain
$D\subset X$ that is closed in $X$, contains the distant coordinate
end, and has compact boundary.  Its intrinsic interior must lie in
$\operatorname{int}X$.  All boundary components and every portion
coincident with $B$ are counted.  Thus $D=X$ is admissible and its cut
is $B$.  Define
\begin{equation}
 a_g(B)=\inf_D\operatorname{Area}_g(\partial_{\mathrm{man}}D).
 \label{eq:g-full-cut-area}
\end{equation}

Attach a compact smooth filling $O$ behind $B$ and extend $g$ smoothly
across $B$, obtaining a connected boundaryless manifold $M$ with
$M\setminus O=\operatorname{int}X$.  Such a filling exists: take a
second copy of a compact truncation of $X$, reverse the collar
direction along $B$, and cap its spherical end boundary with a ball.
This construction does not require an orientation.  The extension
is auxiliary; it will not change any area in
\eqref{eq:g-full-cut-area}.  Since the attached part is compact, the
extended metric can and will be chosen complete.

Let $\mathcal A$ be the bounded finite-perimeter subsets $E$ of $M$
that contain $O$ up to a volume-null set.  Sets agreeing almost
everywhere are identified.  If $\partial^*E$ denotes the reduced
boundary and $\nu_E$ its measure-theoretic outward unit normal, set
\begin{equation}
 P_g(E)=|D\mathbf1_E|_g(M)
       =\mathcal H^{n-1}_g(\partial^*E).
 \label{eq:g-full-perimeter}
\end{equation}
This is perimeter in the extended manifold.  In particular
$P_g(O)=\operatorname{Area}_g(B)$.  Relative perimeter in
$\operatorname{int}X$ would instead assign zero to $O$ and would lose
all contact area.  We use the standard compactness, approximation,
trace, and Gauss--Green results for finite-perimeter sets
\cite{AmbrosioFuscoPallara2000,Maggi2012}.

\begin{lemma}[The full enclosure minimum]
\label{lem:g-enclosure}
Under the preceding hypotheses,
\begin{equation}
 0<a_g(B)=\min_{E\in\mathcal A}P_g(E)
         \le\operatorname{Area}_g(B).
 \label{eq:g-minimum}
\end{equation}
All minimizers are supported in one compact subset of $M$.  The family
of minimizers is compact for $L^1(M)$ convergence, and that convergence
also entails convergence of their full perimeters.

Each minimizer has a closed representative $E$ whose frontier is
$T=\operatorname{supp}|D\mathbf1_E|_g$.  Its exterior $M\setminus E$
is connected and contains the distant end.  There are full smooth
enclosing cuts in $\operatorname{int}X$ with areas tending to $P_g(E)$.
When $3\le n\le7$, these cuts can be chosen to converge to $T$ in
$C^1$.  Every minimizer is outer minimizing: if $F\supset E$ is
bounded and has finite perimeter, then $P_g(F)\ge P_g(E)$.
\end{lemma}

The proof in Section~\ref{subsec:g-proof-enclosure} identifies the
perimeter minimum with the infimum over full smooth cuts. Counting
contact with the original boundary is essential to this identification.

\subsection{Regularity at a smooth obstacle}

The next local statement supplies the regularity used in the enclosure
lemma. It concerns constrained perimeter minimization alone, independently
of global existence or connectedness.

\begin{lemma}[Smooth-obstacle regularity]
\label{lem:g-obstacle}
Let $M^n$ carry a smooth metric, let $O$ have smooth boundary, and let
$E$ locally minimize perimeter among sets containing $O$.
On a relatively compact region meeting $\partial O$, assume that
$\partial O$ has bounded smooth geometry.  Then $E$ has a closed
representative whose topological frontier is
$T=\operatorname{supp}|D\mathbf1_E|_g$.
Its regular part has multiplicity one, and its singular set
$\mathcal S$ satisfies
\[
 \mathcal S=\varnothing\quad(n\le7),\qquad
 \dim_{\mathrm H}\mathcal S\le n-8\quad(n\ge8).
\]
The regular part outside $\partial O$ is smooth and minimal.  Every
contact point in $T\cap\partial O$ is regular, with local graph in
$C^{1,\alpha}\cap W^{2,p}$ for every $0<\alpha<1$ and every
$1<p<\infty$.  In dimensions $3\le n\le7$ the entire frontier is
$C^{1,1}$, with local graph functions in $W^{2,\infty}$.
Moreover, $\mathcal H^{n-1}_g(\mathcal S)=0$ and
$\mathcal H^{n-1}_g(T)=P_g(E)$ wherever the total perimeter is finite.
\end{lemma}

The proof is given in Section~\ref{subsec:g-proof-obstacle}.

\subsection{Detachment from a strict inner barrier}

Contact with the obstacle is possible when its mean curvature is zero.
Strictly negative mean curvature, with the normal pointing toward the
end, instead forces every minimizing enclosure to contain a fixed
exterior collar. The resulting frontier is a locally minimizing
boundary in a smooth ambient neighborhood.

\begin{lemma}[Uniform detachment from a strict inner barrier]
\label{lem:g-detachment}
In the setting of Lemma~\ref{lem:g-enclosure}, suppose that
$H_B=\operatorname{div}_B\nu_B<0$ on every component of $B$, with
$\nu_B$ pointing into $X$.  There is $\varepsilon>0$, depending only
on a fixed collar of $(B,g)$, such that every minimizing filled set
contains the exterior distance collar $0\le t<\varepsilon$ in its
interior.  Consequently its frontier $T$ is compact, nonempty, and
disjoint from $B$, and is locally perimeter minimizing in the smooth
ambient metric.  It is stationary as a whole boundary varifold,
smooth and minimal outside a compact singular set of dimension at
most $n-8$.  If $3\le n\le7$, $T$ is a smooth minimal full cut.

The closed exterior $D=\overline{M\setminus E}$ is connected, has the
same single coordinate end, and is a complete closed subset for the
ambient distance.  Every sequence that is Cauchy for lengths of paths
constrained to $D$ has a limit in $D$ in the same sense, including
when the limit lies on the frontier.
Its frontier is outer
minimizing with all components counted, and
$\mathcal H^{n-1}_g(T)=a_g(B)$.
\end{lemma}

Section~\ref{subsec:g-proof-detachment} proves the collar inclusion and
the completeness assertions. In Section~\ref{sec:rpi:density-repair},
this detached frontier will be the boundary to which we apply the
harmonic-flat Penrose inequality.

Combining the enclosure lemmas with Theorem~\ref{thm:rpi} also gives a
mass bound for a smooth inner boundary that is strictly mean concave
toward the end, without assuming that the boundary itself minimizes area.

\begin{corollary}[The inequality for a strict inner barrier]
\label{cor:g-strict-barrier-penrose}
Let $(X,g)$ be a smooth connected $n$-dimensional exterior, $n\ge3$,
complete including its
nonempty compact smooth boundary $B$, with one Euclidean coordinate
end and compact end complement. Suppose that
\[
 g-\delta=O_2(r^{-\tau}),\qquad \tau>\frac{n-2}{2},\qquad
 R_g\ge0,\quad R_g\in L^1(X,dV_g),\quad
 R_g=O(r^{-n-\beta})\quad(\beta>0).
\]
If $H_B<0$ for the normal pointing into $X$, then
\[
 m_{\rm ADM}(g)\ge\frac12
 \left(\frac{a_g(B)}{\omega}\right)^{(n-2)/(n-1)}.
\]
\end{corollary}

\begin{proof}
Lemma~\ref{lem:g-enclosure} supplies a minimizing enclosure of full
perimeter $a_g(B)$. By Lemma~\ref{lem:g-detachment}, its frontier
is detached from $B$, locally perimeter minimizing, and outer
minimizing; its closed end exterior is connected and complete.
The restricted metric retains the scalar-curvature and end hypotheses,
and is smooth through the frontier. The regularity and area assertions
of the same lemmas therefore permit application of
Theorem~\ref{thm:rpi} to this exterior, giving the displayed bound.
\end{proof}
